\subsection{Legacy L2 under $T_{\mathrm{record}}$}
\label{sec:legacy}

This subsection is a historical oracle baseline under $T_{\mathrm{record}}$, not primary evidence.
Table~\ref{tab:l2legacy} restates the earlier matched chronological comparison that used $T_{\mathrm{record}}$ in $\mathrm{cycle\_norm}$.
Those numbers describe ranking when relative progress to the \emph{recorded} end is known, and when L1 cannot separate from L2.
They are not deleted; they are relabeled as an oracle with a degenerate lifetime definition.

\begin{table}[htbp]
\centering
\caption{Legacy L2 (oracle with $T=T_{\mathrm{record}}$) matched RUL RMSE (cycles; three seeds). Persistence here is persist\_fade on the untruncated chronological tail.}
\label{tab:l2legacy}
\begin{tabular}{@{}lcc@{}}
\toprule
Model & B0005 & B0007 \\
\midrule
persist\_fade & $6.81$ & $30.30$ \\
Linear extrap. & $7.01$ & $28.10$ \\
MLP & $5.22\pm0.99$ & $\mathbf{5.91\pm1.91}$ \\
Patch-TF & $\mathbf{2.28\pm0.96}$ & $9.67\pm2.95$ \\
\bottomrule
\end{tabular}
\end{table}

Dropping $\mathrm{cycle\_norm}$ while keeping the other five channels (L0-drop, same legacy split) collapsed Patch-TF and MLP to about $19$--$22$ RMSE on both cells.
That remains the sharpest single demonstration that the legacy headline error was not a property of capacity/voltage channels alone.

\subsection{Why persist\_fade fails on B0007, and why $T$ is censored}
\label{sec:b0007}

On the legacy split, persist\_fade is 6.81 cycles on B0005 and 30.3 on B0007.
Fig.~\ref{fig:soh}B shows the mechanism: B0007's recorded SOH ends near 0.76 and never crosses 0.7, so $T$ equals the last observed cycle, and late increments are often positive.
A window fade rate near zero maps to an enormous implied RUL.
The failure therefore has two coupled causes: the local fade is non-monotone, and the target is \emph{right-censored}, so extrapolating to 0.7 asks the baseline to invent a decline that was never observed.
The L0-valid recursive walk, initialized by persist\_fade then counting down by one cycle, is $15.53$ on B0007 versus $6.14$ on B0005 (Table~\ref{tab:persist}).
persist\_hold remains $1.0$ because it copies true $\mathrm{RUL}_{t-1}$ and is an oracle.
“Learned models beat persistence under L0 on B0007” is therefore true for persist\_fade ($30.29$) and must not be written against persist\_hold.
Because B0007 is censored, its absolute RUL RMSE is a record-horizon error and should not be compared numerically with B0005/B0006; it is used here as a stress case for the L1P mechanism and for persist\_fade, not as a lifetime benchmark.

\section{Discussion}
\label{sec:discussion}

\paragraph{Finding 1: information level changes the ranking.}
Under Severson Protocol~B ($n{=}24$), median MLP RMSE falls from $222.1$ at L0 to $12.0$ at L2, and L0 exceeds L2 on every scored cell (Fig.~\ref{fig:seversonb}).
That leave-target-cell-out ladder is the primary support for $\operatorname{Rank}(\cdot\mid\mathcal{I}_t)=f(\mathcal{I}_t)$.
The L0 error is in the same regime as persist\_fade ($242.8$), so the ladder is not an artefact of a weak model family.

\paragraph{Finding 2: the L1P inversion is a support-and-learnability effect.}
NASA Protocol~A adds a narrower second lesson, and the mechanism analysis sharpens it.
L1P can beat L2 even when SOH matching distances stay small (Fig.~\ref{fig:l1pnasa}), because the L2 clock is evaluated outside its training support and the fitted network compresses late-life variation, while L1P is an SOH-indexed local average that remains inside its matching neighbourhood (Section~\ref{sec:mechanism}).
This is not a case of “less information is better”.
It is a case in which additional information was \emph{sufficient but not learnable} on the scored windows under the chosen fit protocol, a distinction that a single RMSE column cannot reveal.
The inversion disappears under Protocol~B (L1P $30.3$ vs.\ L2 $12.0$) and reverses on CALCE CS2\_36 (L1P $60.69$ vs.\ L2 $2.20$), so the correct claim is conditional, not universal.

\paragraph{Finding 3: Protocol~A and Protocol~B answer different questions.}
Equation~\eqref{eq:affine} explains why L1 and L2 look alike under Protocol~A: they are affine rescalings of the same~$t$ (Table~\ref{tab:affine}).
Scale-control runs confirm empirical robustness of the trained networks, but they do not prove invariance.
Protocol~B breaks the practical L1$\approx$L2 pattern because a single sibling prior cannot restore target-specific relative life across $T\in[147,1933]$.
Under Protocol~A, raw~$t$ can also beat L2 when the late-life test set is essentially a within-cell countdown.
That pattern is a Protocol~A artifact; it does not overturn the Protocol~B result, and it is the reason Section~\ref{sec:ladder} is labelled exploratory.

\paragraph{Relation to benchmark platforms.}
BatteryML~\cite{zhang2024batteryml} and BatteryLife~\cite{tan2025batterylife} make data pipelines and model comparisons reproducible; this paper addresses a different axis, the auditability of the information condition.
The two are complementary: dataset and model cards can carry an “information level / fit protocol / $T$ definition / scaler source / support status” row for every reported score, and the ladder can be implemented as a wrapper around their evaluation harness without changing the models.
We do not claim that current platforms are leaky; we claim that their reporting cards leave the information condition implicit, which is enough for two incomparable scores to appear side by side.

\paragraph{Practical guidance.}
For an online BMS claim, L2 must be reported as a named retrospective oracle and never as a deployable result.
When $T$ is unknown and a comparable fleet exists, an SOH-indexed sibling prior is a reasonable online-valid plug-in, with the caveat that it needs SOH coverage and a stable local SOH--RUL relation; when no fleet exists, keep the L0 recursive walk and report the oracle only as an upper bound.
For a benchmark claim, report Protocol~B; for a within-cell retrospective study, report Protocol~A and state that L1/L1P may be undefined without siblings.

\paragraph{Limitations.}
NASA and CALCE stay Protocol~A diagnostics, and the mechanism analysis rests on three NASA cells.
Severson Protocol~B uses a stratified 24-cell stress subset of 119 usable cells, not a random population sample.
B0007 is right-censored, so its absolute errors are record-horizon errors.
Beyond NASA cycle aggregates, the extracts omit full in-cycle V/I/T trajectories.
All scores are cycle-level after a completed discharge.
The SOH head is a diagnostic, not a validated SOH estimator, and its repair is demonstrated only on NASA.
Finally, L1P is a local prior: it is only as good as the sibling coverage, and its failure on CALCE CS2\_36 is part of the result rather than an outlier to be removed.

\section{Conclusion}
\label{sec:conclusion}

RUL benchmarks should encode prediction-time information in the protocol itself.
Information-consistent evaluation couples four things that are usually reported separately: a named information level, a declared fit protocol, a provenance and support audit, and a fillable report row.
Under first-EOL labels and named L0/L1/L1P/L2 cycle channels, leave-target evaluation on Severson ($n{=}24$) shows that information level changes RUL error (L0 $222.1$ vs.\ L2 $12.0$; $p=1.2{\times}10^{-7}$).
NASA Protocol~A further shows a setting-specific reversal, L1P beating L2, which we trace to the oracle clock leaving its training support while an SOH-indexed local prior stays inside its matching neighbourhood.
Future tables should report information level, fit protocol, the $T$ definition and censoring status, channel support, and persist\_fade with RMSE, not the architecture name alone.

\section*{Appendix. OLS identity check}
% OLS identity check (appendix only; not a predictive baseline)
\begin{table}[htbp]
\centering
\caption{Appendix OLS identity check on evaluated chronological windows (RUL$>0$).
OLS is an algebraic diagnostic of $\mathrm{RUL}=\max(0,T-t)$, not an out-of-sample predictive baseline: it is fit with an intercept on the same scored windows.
Under that identity, OLS RMSE on raw~$t$ is numerically zero on NASA and CALCE.
$\rho$ is Pearson correlation; $\rho(\mathrm{L1},\mathrm{L2})=1$ because L1 and L2 are affine in~$t$.
L1P is not affine in~$t$.}
\label{tab:clock}
\begin{tabular}{@{}llcccc@{}}
\toprule
Corpus & Cell & $\rho(\mathrm{L1},\mathrm{L2})$ & $\rho(\mathrm{L1P},\mathrm{L2})$ & OLS $t$ & OLS L1P \\
\midrule
nasa & B0005 & 1.000 & -0.703 & 0 & 14.20 \\
nasa & B0006 & 1.000 & -0.598 & 0 & 15.57 \\
nasa & B0007 & 1.000 & -0.824 & 0 & 13.28 \\
calce & CS2\_35 & 1.000 & -0.122 & 0 & 98.18 \\
calce & CS2\_36 & 1.000 & -0.706 & 0 & 101.04 \\
calce & CS2\_37 & 1.000 & -0.654 & 0 & 99.90 \\
calce & CS2\_38 & 1.000 & -0.621 & 0 & 96.29 \\
\bottomrule
\end{tabular}
\end{table}


\section*{Appendix. Seed-level contrasts and full SOH table}
% auto-generated appendix tables (moved from generated_tables.tex)
\begin{table}[htbp]
\centering
\caption{SOH RMSE (dimensionless; mean$\pm$s.d.\ over three seeds) under the same information ladder.}
\label{tab:soh}
\begin{tabular}{@{}llccc@{}}
\toprule
Model & Cell & L2 & L1 & L0 \\
\midrule
mlp & B0005 & 0.7063$\pm$0.1843 & 0.7007$\pm$0.1826 & 0.3302$\pm$0.0660 \\
mlp & B0006 & 0.8120$\pm$0.2939 & 0.8180$\pm$0.2988 & 0.2879$\pm$0.0462 \\
mlp & B0007 & 0.5625$\pm$0.1420 & 0.5640$\pm$0.1332 & 0.3358$\pm$0.0264 \\
patch-tf & B0005 & 0.3593$\pm$0.0567 & 0.3583$\pm$0.0565 & 0.5132$\pm$0.1256 \\
patch-tf & B0006 & 0.5564$\pm$0.3209 & 0.5548$\pm$0.3199 & 0.4427$\pm$0.0678 \\
patch-tf & B0007 & 0.4028$\pm$0.0108 & 0.4044$\pm$0.0119 & 0.6538$\pm$0.4076 \\
\bottomrule
\end{tabular}
\end{table}

\begin{table}[htbp]
\centering
\caption{Seed-level paired contrasts of RUL RMSE (selected). The experimental unit is the random seed on a fixed cell ($n=3$), not the cell. Wilcoxon $p{=}0.25$ is the smallest two-sided value when all three seeds move in the same direction. Full contrasts are in \texttt{outputs/info\_levels\_v2/stats.json}.}
\label{tab:stats}
\begin{tabular}{@{}llllccc@{}}
\toprule
Contrast & Corpus & Model & Cell & $\Delta$RMSE & Cohen's $d$ & $p$ \\
\midrule
L2-vs-L0 & calce & mlp & CS2\_35 & 30.69 & 13.98 & 0.250 \\
L2-vs-L1 & calce & mlp & CS2\_35 & 0.27 & 0.15 & 0.750 \\
L2-vs-L0 & calce & mlp & CS2\_36 & 38.13 & 21.52 & 0.250 \\
L2-vs-L1 & calce & mlp & CS2\_36 & 0.08 & 0.07 & 0.750 \\
L2-vs-L0 & calce & mlp & CS2\_37 & 42.55 & 22.73 & 0.250 \\
L2-vs-L1 & calce & mlp & CS2\_37 & 0.14 & 0.13 & 0.500 \\
L2-vs-L0 & calce & mlp & CS2\_38 & 34.13 & 10.86 & 0.250 \\
L2-vs-L1 & calce & mlp & CS2\_38 & 0.53 & 0.30 & 0.500 \\
L2-vs-L0 & nasa & mlp & B0005 & -11.81 & -4.33 & 0.250 \\
L2-vs-L1 & nasa & mlp & B0005 & 0.04 & 0.01 & 0.750 \\
L2-vs-L0 & nasa & mlp & B0006 & -4.16 & -3.00 & 0.250 \\
L2-vs-L1 & nasa & mlp & B0006 & -0.10 & -0.06 & 0.250 \\
L2-vs-L0 & nasa & mlp & B0007 & -8.97 & -2.73 & 0.250 \\
L2-vs-L1 & nasa & mlp & B0007 & -0.03 & -0.01 & 1.000 \\
L2-vs-L0 & nasa & patch-tf & B0005 & -8.95 & -4.67 & 0.250 \\
L2-vs-L1 & nasa & patch-tf & B0005 & 0.01 & 0.00 & 1.000 \\
L2-vs-L0 & nasa & patch-tf & B0006 & -5.49 & -19.33 & 0.250 \\
L2-vs-L1 & nasa & patch-tf & B0006 & -0.01 & -0.03 & 0.750 \\
L2-vs-L0 & nasa & patch-tf & B0007 & -4.41 & -2.03 & 0.250 \\
L2-vs-L1 & nasa & patch-tf & B0007 & 0.02 & 0.01 & 0.500 \\
\bottomrule
\end{tabular}
\end{table}



\section*{Appendix. Reporting template}

\begin{table}[htbp]
\centering
\footnotesize
\setlength{\tabcolsep}{3pt}
\caption{Minimum report row (copy per model $\times$ cell $\times$ level).}
\label{tab:template}
\begin{tabular}{@{}p{0.14\linewidth}p{0.07\linewidth}p{0.15\linewidth}p{0.08\linewidth}p{0.08\linewidth}p{0.10\linewidth}p{0.09\linewidth}@{}}
\toprule
Model & Level & Cycle channel & $W$, $H$ & Scaler & Support & RUL RMSE \\
\midrule
Patch-TF & L0 & raw $t$ & 50, 5 & train & in & \ldots \\
Patch-TF & L1 & $t/T_{\mathrm{med}}$ & 50, 5 & train & in & \ldots \\
Patch-TF & L2 & $t/T_{\mathrm{target}}$ & 50, 5 & train & out (oracle) & \ldots \\
persist\_fade & --- & SOH only & 50, --- & --- & --- & \ldots \\
\bottomrule
\end{tabular}
\end{table}

Required fields: corpus, cell id, $T$ definition and whether EOL was reached, whether post-EOL windows are dropped, persist\_fade RMSE, fit protocol, channel support status, and the information level of every learned row.
Optional: SOH RMSE as an auxiliary control; PICP/PINAW in the appendix.

\section*{Appendix. Uncertainty (not a calibration claim)}
% auto-generated UQ appendix; do not edit by hand
\begin{table}[htbp]
\centering
\caption{Appendix. MC-Dropout 95\% intervals ($n_{\mathrm{MC}}=20$) for Patch-TF RUL on NASA. PICP is empirical coverage of mean$\pm 1.96$ std; PINAW is width normalized by the test RUL range. Not used as a calibration claim.}
\label{tab:uq}
\begin{tabular}{@{}lcccc@{}}
\toprule
Cell & L2 PICP & L2 PINAW & L0 PICP & L0 PINAW \\
\midrule
B0005 & 0.058$\pm$0.025 & 0.599$\pm$0.077 & 0.522$\pm$0.174 & 0.879$\pm$0.167 \\
B0006 & 0.056$\pm$0.048 & 0.709$\pm$0.041 & 0.528$\pm$0.210 & 0.878$\pm$0.265 \\
B0007 & 0.217$\pm$0.075 & 0.607$\pm$0.097 & 0.377$\pm$0.280 & 0.792$\pm$0.165 \\
\bottomrule
\end{tabular}
\end{table}


\section*{Statements and Declarations}
\paragraph{Funding.} None.
\paragraph{Competing interests.} None.
\paragraph{Data availability.}
NASA PCoE cells B0005--B0007~\cite{saha2007nasa}; CALCE CS2 cells~\cite{he2011calce}; Severson/MATR batches~\cite{severson2019nature}.
Replication scripts and adopted experiment outputs: \path{scripts/run_info_levels_v2.py}, \path{scripts/run_revision_controls.py}, \path{scripts/run_l1p_mechanism.py}, \path{experiments/info_consistent/}, \path{outputs/info_levels_v2/}, \path{outputs/info_levels_severson_loco/}, \path{outputs/revision_controls/}, \path{outputs/l1p_mechanism/}.
A versioned code bundle is maintained in the project git repository; the public GitHub URL and Zenodo DOI will be inserted here upon first deposit (see \path{CODE_AVAILABILITY.md}).

\begin{thebibliography}{99}

%% ---- foundational corpora / protocol anchors (pre-2025 minority) ----
\bibitem{saha2007nasa}
B.~Saha, K.~Goebel,
Battery Data Set,
NASA Ames Prognostics Data Repository,
Moffett Field, CA (2007).

\bibitem{saha2009compare}
B.~Saha, K.~Goebel, J.~Christophersen,
Comparison of prognostic algorithms for estimating remaining useful life of batteries,
Trans.\ Inst.\ Meas.\ Control 31~(3--4) (2009) 293--308.

\bibitem{he2011calce}
W.~He, N.~Williard, M.~Osterman, M.~Pecht,
Prognostics of lithium-ion batteries based on Dempster--Shafer theory and the Bayesian Monte Carlo method,
J.\ Power Sources 196~(23) (2011) 10314--10321.

\bibitem{kaufman2012leakage}
S.~Kaufman, S.~Rosset, C.~Perlich, O.~Stitelman,
Leakage in data mining: Formulation, detection, and avoidance,
ACM Trans.\ Knowl.\ Discov.\ Data 6~(4) (2012) 1--21.

\bibitem{gal2016}
Y.~Gal, Z.~Ghahramani,
Dropout as a Bayesian approximation: representing model uncertainty in deep learning,
in: Proc.\ 33rd Int.\ Conf.\ Mach.\ Learn., PMLR 48 (2016) 1050--1059.

\bibitem{vaswani2017}
A.~Vaswani, N.~Shazeer, N.~Parmar, J.~Uszkoreit, L.~Jones, A.N.~Gomez, L.~Kaiser, I.~Polosukhin,
Attention is all you need,
in: Adv.\ Neural Inf.\ Process.\ Syst.\ 30 (2017) 5998--6008.

\bibitem{birkl2017}
C.R.~Birkl, M.R.~Roberts, E.~McTurk, P.G.~Bruce, D.A.~Howey,
Degradation diagnostics for lithium ion cells,
J.\ Power Sources 341 (2017) 373--386.

\bibitem{xiong2018soh}
R.~Xiong, L.~Li, J.~Tian,
Towards a smarter battery management system: a critical review on battery state of health monitoring methods,
J.\ Power Sources 405 (2018) 18--29.

\bibitem{lipu2018}
M.S.H.~Lipu, M.A.~Hannan, A.~Hussain, M.M.~Hoque, P.J.~Ker, S.M.A.~Saad, A.~Ayob,
A review of state of health and remaining useful life estimation methods for lithium-ion battery in electric vehicles: challenges and recommendations,
J.\ Clean.\ Prod.\ 205 (2018) 115--133.

\bibitem{hu2019rser}
X.~Hu, F.~Feng, K.~Liu, L.~Zhang, J.~Xie, B.~Liu,
State estimation for advanced battery management: key challenges and future trends,
Renew.\ Sustain.\ Energy Rev.\ 114 (2019) 109334.

\bibitem{severson2019nature}
K.A.~Severson, P.M.~Attia, N.~Jin, N.~Perkins, B.~Jiang, Z.~Yang, M.H.~Chen, M.~Aykol, P.K.~Herring, D.~Fraggedakis, M.Z.~Bazant, S.J.~Harris, W.C.~Chueh, R.D.~Braatz,
Data-driven prediction of battery cycle life before capacity degradation,
Nat.\ Energy 4 (2019) 383--391.

\bibitem{roman2021}
D.~Roman, S.~Saxena, V.~Robu, M.~Pecht, D.~Flynn,
Machine learning pipeline for battery state-of-health estimation,
Nat.\ Mach.\ Intell.\ 3 (2021) 447--456.

\bibitem{sulzer2021joule}
V.~Sulzer, P.~Mohtat, A.~Aitio, S.~Lee, Y.T.~Yeh, F.~Steinbacher, M.U.~Khan, J.W.~Lee, J.B.~Siegel, A.G.~Stefanopoulou, D.A.~Howey,
The challenge and opportunity of battery lifetime prediction from field data,
Joule 5 (2021) 1934--1955.

\bibitem{geslin2023joule}
A.~Geslin, B.~van Vlijmen, X.~Cui, A.~Bhargava, P.A.~Asinger, R.D.~Braatz, W.C.~Chueh,
Selecting the appropriate features in battery lifetime predictions,
Joule 7 (2023) 1956--1965.

\bibitem{kapoor2023leakage}
S.~Kapoor, A.~Narayanan,
Leakage and the reproducibility crisis in machine-learning-based science,
Patterns 4~(9) (2023) 100804.

\bibitem{nie2023patchtst}
Y.~Nie, N.H.~Nguyen, P.~Sinthong, J.~Kalagnanam,
A time series is worth 64 words: long-term forecasting with transformers,
in: Int.\ Conf.\ Learn.\ Represent.\ (2023).

\bibitem{rosenblatt2024}
M.~Rosenblatt, L.~Tejavibulya, R.~Jiang, S.M.~Noble, D.~Scheinost,
Data leakage inflates prediction performance in connectome-based machine learning models,
Nat.\ Commun.\ 15 (2024) 1829.

\bibitem{zhang2024batteryml}
H.~Zhang, X.~Gui, S.~Zheng, Z.~Lu, Y.~Li, J.~Bian,
BatteryML: An open-source platform for machine learning on battery degradation,
in: Int.\ Conf.\ Learn.\ Represent.\ (2024).

\bibitem{li2024lifetime}
T.~Li, Z.~Zhou, A.~Thelen, D.A.~Howey, C.~Hu,
Predicting battery lifetime under varying usage conditions from early aging data,
Cell Rep.\ Phys.\ Sci.\ 5~(4) (2024) 101891.

%% ---- 2025--2026 majority ----
\bibitem{wang2024gce}
S.~Wang, R.~Zhou, Y.~Ren, M.~Jiao, H.~Liu, C.~Lian,
Advanced data-driven techniques in AI for predicting lithium-ion battery remaining useful life: a comprehensive review,
Green Chem.\ Eng.\ 6~(2) (2025) 139--153.

\bibitem{zhang2025batlinet}
H.~Zhang, Y.~Li, S.~Zheng, Z.~Lu, X.~Gui, W.~Xu, J.~Bian,
Battery lifetime prediction across diverse ageing conditions with inter-cell deep learning,
Nat.\ Mach.\ Intell.\ 7 (2025) 270--277.

\bibitem{tan2025batterylife}
R.~Tan, W.~Hong, J.~Tang, X.~Lu, R.~Ma, X.~Zheng, J.~Li, J.~Huang, T.-Y.~Zhang,
BatteryLife: A comprehensive dataset and benchmark for battery life prediction,
arXiv:2502.18807 (2025).

\bibitem{madani2025groupkfold}
S.S.~Madani, M.~H\'{e}bert, L.~Boulon, A.~Lupien-B\'{e}dard, F.~Allard,
Comparative analysis of ML and DL models for data-driven SOH estimation of LIBs under diverse temperature and load conditions,
Batteries 11~(11) (2025) 393.

\bibitem{xu2025meas}
Z.~Xu, J.~Mao, Y.~Han, H.~Yuan,
A comprehensive review of lithium-ion battery remaining useful life prediction: methodologies, datasets, performance metrics, and future perspectives,
Meas.\ Sci.\ Technol.\ 36~(8) (2025) 082001.

\bibitem{yang2025batteries}
Y.~Yang, S.~Wang, Y.~Zhou, C.~Fernandez,
A critical review of AI-based battery remaining useful life prediction for energy storage systems,
Batteries 11~(10) (2025) 376.

\bibitem{deng2025piml}
Y.~Deng, X.~Le, T.T.~Nguyen, C.~Gogu,
A generic physics-informed machine learning framework for battery remaining useful life prediction using small early-stage lifecycle data,
Appl.\ Energy 384 (2025) 125314.

\bibitem{oyewole2025uq}
I.~Oyewole, W.~Hassanieh, M.~Chelbi, A.~Chehade,
Uncertainty-aware deep learning with physics-informed bayesian sampling for lithium-ion battery prognostics,
Appl.\ Energy 402 (2025) 126880.

\bibitem{aliakbari2025fusedr}
S.~Aliakbari, M.~Masih-Tehrani,
FuseDR-Net: a fusion-based network for battery remaining useful life prediction using causally decomposed health indicators,
Energy Convers.\ Manag.\ X 28 (2025) 101341.

\bibitem{ge2025koopman}
X.~Ge, Y.~Jiang, B.~Xu,
Deep Koopman operator-based remaining useful life prediction of Lithium-ion batteries under multi-condition scenarios,
J.\ Energy Storage 119 (2025) 116369.

\bibitem{li2025adapt}
J.~Li, Z.~Wang, D.~Tian, Z.~Sun, Y.~Niu,
A novel RUL prediction framework based on the adaptability feature perception fusion model method,
J.\ Energy Storage 119 (2025) 116322.

\bibitem{huang2025rulmamba}
J.~Huang, L.~Liu, H.~Zhao, T.~Li, B.~Li,
RUL-Mamba: Mamba-based remaining useful life prediction for lithium-ion batteries,
J.\ Energy Storage 120 (2025) 116376.

\bibitem{tian2025energy}
A.~Tian, H.~Yu, Z.~Hu, Y.~Wang, T.~Wu, J.~Jiang,
A novel remaining useful life prediction method based on CNN-Attention combined with SMA-GPR,
Energy 321 (2025) 135233.

\bibitem{yi2025jps}
X.~Yi, J.~Xiao, G.~Lei, X.-H.~Hu, Z.~Feng,
A lithium-ion battery remaining useful life prediction model based on multilayer perceptron expert networks and temporal feature composition,
J.\ Power Sources 659 (2025) 238371.

\bibitem{wang2025jes}
K.~Wang, J.~Wang, Z.~Wang, H.~Li,
An enhanced deep learning framework for state of health and remaining useful life prediction of lithium-ion battery based on discharge fragments,
J.\ Energy Storage 107 (2025) 114952.

\bibitem{ali2026loco}
B.~Ali, I.~Matloob, N.~Alkhalifa, A.K.~Alkhalifa, N.~Asim, S.~Kanwal, S.~Inam,
Deep learning-based RUL and SOH prediction of lithium-ion batteries using LOCO,
Sci.\ Rep.\ 16 (2026) 27853.

\bibitem{barua2026predictor}
A.~Barua, J.M.~Collins, M.R.~Islam, S.~Barua,
Predictor structure modulates validation inflation and cell-transfer reliability in battery state-of-health estimation,
Batteries 12~(9) (2026) 342.

\bibitem{sun2026sgdits}
Y.~Sun, F.~Gao,
SG-DiTs: A physics-informed diffusion-transformer hybrid framework for reliable battery RUL prediction,
J.\ Energy Storage 152 (2026) 120479.

\bibitem{fan2026meas}
Y.~Fan, Z.~Xu, H.~Zhu, Y.~Yu,
Physics-informed signal decomposition and data-driven RUL estimation for lithium-ion batteries,
Meas.\ Sci.\ Technol.\ 37~(18) (2026) 186103.

\bibitem{li2026dpmamba}
X.~Li, H.~Zhang,
DpMamba: a dual-path state space model for reliable remaining useful life prediction of lithium-ion batteries,
Meas.\ Sci.\ Technol.\ 37~(34) (2026) 346105.

\end{thebibliography}


\end{document}